\subsection{Lower bounds and Communication Complexity}
As we gain and store every more information, we constantly need more effective
ways to process and store that information. To do this more effectively, we can
either create more powerful hardware, or create ever more efficient algorithms.
We will consider the latter course. Many algorithm designers have created
enormously efficient algorithms for solving many core problems in querying and
optimization, but improvement does beg the question--how fast can we make these
algorithms? Certainly, asked to check if two numbers are the same, our
natural inclination might be to simply compare all of the digits, resulting in
a protocol running in $\mathcal{O}(n)$ time. However, is it possible to do
better than this? And can we prove it? While classical computing might provide
tools for the case of equality---No, we cannot do better---many more complex
questions such as optimization using linear programs cannot be proven using
these classical tools. Communication Complexity provides a model under which we
can transform these logical processes, for which we lack adequate mathematical
tools to prove bounds on, into combinatorial ones for which we can find
solutions much more easily.

%%% TODO: We should decide if we want to keep this section or not. I don't think we need it, since we are able to motivate our problem using algorithm development above, but whether we keep it or not I dont think it fits well in this section
\subsection{Motivation from multi-party traversal case studies}
Large-scale transportation and infrastructure disruptions create a natural model for multi-party graph traversal. Consider the 2010 European volcanic ash shutdown: air-navigation service providers, airports, and airlines each control disjoint portions of airspace and schedules. Rerouting passengers is a shortest-path problem in a layered graph whose feasible vertices and edges change over time. No party can broadcast its full real-time state, both for operational and privacy reasons, so route feasibility must be computed with limited information exchange.

Similar graph-traversal constraints appear in urban transit shutdowns and phased restorations, where distinct agencies own stations and tunnels. After Hurricane Sandy, service availability evolved over weeks, and agencies coordinated alternate routes under partial knowledge of what infrastructure was open. Highway detours after the I-95 bridge collapse in Philadelphia illustrate the same phenomenon in a road network: different operators own highway segments and toll roads, and a single failed cut edge forces traffic onto alternative paths with asymmetric knowledge of capacities.

The maritime and logistics examples are equally instructive. The 2021 Suez Canal blockage forced carriers to reroute around the Cape of Good Hope or wait, a decision that depends on schedule and capacity information held privately by shipping lines and port operators. When the Baltimore Key Bridge collapsed, access to the port was blocked and cargo was diverted to alternate ports, again requiring coordination across a sea--port--land graph with private operational constraints.

Communication bottlenecks become even sharper in critical infrastructure networks. During the 2021 Texas ERCOT blackout, grid operators and utilities effectively solve a constrained flow and connectivity problem as generation and transmission nodes go offline. The Colonial Pipeline cyber disruption similarly forced fuel logistics to reroute through alternate depots and terminals with incomplete visibility. Finally, the 2021 Meta backbone withdrawal and the 2016 Dyn DNS attack show that digital services rely on reachability in routing and dependency graphs, where policy and security constraints limit information sharing across administrative domains.

These cases motivate a theoretical study of multi-party graph traversal: agents must collaboratively decide reachability and shortest paths while their private subgraphs, capacities, or weights are not fully shareable. The fundamental question is how much communication is necessary, and how randomized protocols, sketches, or partial disclosures can reduce the burden while preserving correctness.



\subsection{Why communication is the bottleneck}
In the model we consider, each party has unlimited local computation but limited ability to share data. This is realistic: air traffic control centers, network operators, or utilities can compute detailed local routes, but cannot afford to broadcast all local states and constraints due to bandwidth, privacy, regulatory, or security limitations. The communication cost, not the local computation, is what constrains real-time coordination.

Communication complexity makes this separation explicit. In particular, it shows that even when each party can compute arbitrarily powerful local summaries, some global questions still require a large number of bits to be exchanged. For example, deciding whether two privately held sets intersect already requires linear communication in the worst case. This lower bound appears again in many traversal problems via reductions that embed set disjointness or equality into graph reachability. Thus, communication is not simply a practical nuisance: it is an inherent barrier that must be designed around.


In this report, we focus on studying the specific case of graph problems, in
which two or more parities are given a subset of the graph and asked to compute
its properties. Primarily we research vertex reachability, that is, given two
vertices $(u,v)$, identify it there is a path between them which uses only the
subsets of the graph owned by the communicating parties.
